\documentclass{article}

\usepackage[preprint]{neurips_2026}
\usepackage[T1]{fontenc}
\usepackage[utf8]{inputenc}
\usepackage{hyperref}
\usepackage{url}
\usepackage{booktabs}
\usepackage{amsfonts}
\usepackage{nicefrac}
\usepackage{microtype}
\usepackage{xcolor}
\usepackage{amsmath}
\usepackage{amssymb}
\usepackage{graphicx}

\title{Shift-Factored Predictive Transport for Probabilistic Regression}

\author{%
  Author names withheld for branch development
}

\begin{document}

\maketitle

\begin{abstract}
We study probabilistic regression under deployment shift, where source-trained
predictive laws become biased, under-dispersed, and poorly calibrated on target
data. We propose Shift-Factored Predictive Transport (SPT-Reg), a test-time
framework that adapts predictive distributions by composing output-space prior
transport, predictive-subspace alignment, and representation-aware uncertainty
inflation. The framework is predictor-family agnostic through a shared
predictive container spanning Gaussian, quantile, histogram, and ordered-bin
regression outputs. Conformal wrapping restores finite-sample coverage, while
prediction-powered inference enables valid target-domain confidence intervals
with limited labels. We show how this factorized transport view improves
calibration, tail coverage, selective prediction, and inferential efficiency
under realistic regression shifts without requiring target retraining.
\end{abstract}

\section{Introduction}

Many regression deployment failures are not caused by a complete collapse of
$p(y \mid x)$, but by target-domain distortions in predictive calibration,
output prevalence, and uncertainty scale. This paper develops the hypothesis
that these failures can often be repaired by transporting the source predictive
law rather than retraining the regressor from scratch.

The branch implementation targets four predictive families already exposed by
\texttt{torchregress}: Gaussian, quantile, bar/histogram, and ordered-bin
predictors. The paper focus is calibration, coverage, selective prediction, and
target-domain inference under limited labels, rather than raw RMSE alone.

Our empirical position is deliberately strong: SPT-Reg should be evaluated as a
competing method against target fine-tuning, raw conformalization, feature-only
normalization, and family-matched BinnedPDF baselines, not as an isolated
ablation stack. This framing matters for NeurIPS review because it turns the
paper from a library integration story into a concrete claim about
label-efficient target adaptation for predictive laws.

\section{Method}

SPT-Reg composes five stages:
\begin{enumerate}
  \item output-space prior transport on a shared support grid,
  \item significance-weighted predictive-subspace alignment,
  \item representation-aware uncertainty inflation,
  \item optional conformal wrapping, and
  \item prediction-powered inference on adapted predictions.
\end{enumerate}

The current branch implementation exposes this flow through a reusable
\texttt{ShiftFactoredPredictiveTransport} orchestrator in
\texttt{torchregress.test\_time}.

Internally, the method maps each predictive family to a common
\texttt{PredictiveBatch} representation and then to a shared support-grid
density. This lets the same transport operator act on Gaussian, quantile,
histogram, and ordered-bin / BinnedPDF outputs without changing the surrounding
evaluation contract. In the current implementation, prior transport is driven
by confidence-filtered and locally consistent target rows, significance-weighted
subspace alignment is only applied when the predictor can be rerun on aligned
inputs, and representation-shift calibration inflates either Gaussian variance
or density temperature depending on the predictive family.

The branch also treats conformal wrapping and prediction-powered inference as
first-class method stages rather than post-processing afterthoughts. This is
important for the paper narrative: transport is used to improve sharpness and
residual quality under shift, while conformal prediction and PPI restore
finite-sample validity for intervals and downstream estimands.

\section{Experiments}

The experimental program is organized into three tracks:
\begin{enumerate}
  \item synthetic factorized-shift benchmarks,
  \item real tabular competing-method evaluations on both a small built-in
  dataset and a larger YearPredictionMSD-style dataset, and
  \item an optional photo-$z$ compatibility track kept lightweight in
  \texttt{torchregress} because the main photo-$z$ effort has moved to
  \texttt{torchz}.
\end{enumerate}

All main tables and figures should be generated from machine-readable summaries
under \texttt{reports/neurips\_spt\_reg/}.

\paragraph{Competing methods.}
The main comparison set includes raw source predictors, feature-stat
normalization, significant-subspace alignment alone, prior transport alone,
raw split conformal, small-label target refitting, and the full SPT-Reg
pipeline. We report each of these on Gaussian predictors where applicable, and
we additionally compare ordered-bin / BinnedPDF predictors with and without
SPT-Reg in the synthetic and tabular tracks. In \texttt{torchregress}, the
photo-$z$ track is treated as an optional compatibility benchmark rather than
the core submission path.

\paragraph{Metrics.}
The primary tables emphasize probabilistic and decision-facing quantities:
negative log-likelihood, CRPS, coverage at $90\%$, interval width,
risk-coverage summaries, and target-domain PPI interval width and coverage.
For the photo-$z$ benchmark we additionally report NMAD, catastrophic outlier
rate, and high-redshift error slices so that improvements in predictive
transport do not hide tail failures that matter in scientific deployment, but
those diagnostics should be considered supporting evidence here unless the
final photo-$z$ evaluation is pulled in from \texttt{torchz}.

\paragraph{Artifact generation.}
The branch includes a single render entrypoint,
\texttt{tools/render\_spt\_reg\_paper\_artifacts.py}, which writes the
synthetic, small-tabular, and larger YearPredictionMSD-style summary JSON
files by default, plus an artifact manifest used by the paper reproducibility
notes. An optional flag adds the local photo-$z$ compatibility benchmark. This
keeps the manuscript tied to committed outputs rather than ad hoc notebook
tables while avoiding a hard dependency on the older photo-$z$ surface in this
repository.

\section{Limitations}

SPT-Reg is expected to degrade when the target shift is primarily
non-transportable, such as severe conditional mechanism drift outside the
learned predictive subspace. The paper should report these failure modes
explicitly and separate validity claims from efficiency claims.

The discretized prior-transport stage is also an approximation to continuous
output shift, not a claim that regression label shift is solved in full
generality. The paper therefore needs to show when transport helps, when it
merely supports valid conformal or PPI correction, and when small-label target
refitting remains the stronger option.

\bibliographystyle{plainnat}
\bibliography{refs}

\appendix

\section{Appendix}

This appendix will collect:
\begin{itemize}
  \item theorem and proposition proof sketches,
  \item full competing-method tables,
  \item ablations for each transport stage,
  \item dataset split details,
  \item additional photo-$z$ diagnostics, and
  \item a reproducibility checklist keyed to
  \texttt{reports/neurips\_spt\_reg/artifact\_manifest\_latest.json}.
\end{itemize}

\end{document}
